\clearpage
\section{Kiểm định Hotelling \texorpdfstring{$T^2$}{T²}: trung bình có khác giả thuyết?}
\label{sec:hotelling}

\subsection{Vì sao không kiểm định riêng từng biến?}
Giả sử cần kiểm tra $H_0:\vmu=(\mu_{01},\mu_{02})'$. Làm hai kiểm định
$t$ mức 5\% không tạo ra một kiểm định chung mức 5\%: xác suất có ít nhất
một bác bỏ sai có thể lớn hơn 5\%. Hơn nữa, một hướng như $X_1-X_2$ có
thể ít biến thiên dù từng biến riêng có phương sai lớn. Sự khác biệt theo
hướng đó có thể bị bỏ qua nếu chỉ kiểm định từng thành phần.

Hotelling $T^2$ dùng \emph{toàn bộ} ma trận hiệp phương sai để đánh giá
vectơ chênh lệch giữa trung bình mẫu và giả thuyết. Nó là phép mở rộng
trực tiếp của
\[
 t^2=n(\bar x-\mu_0)^2/s^2.
\]
\nguon{Chương 5, §5.2, tr.~210--212, công thức (5-1)--(5-4)}

\subsection{Thống kê và phân phối chính xác}
Ta xét giả thuyết hai phía
\[
 H_0:\vmu=\vmu_0,\qquad H_1:\vmu\ne\vmu_0.
\]
\begin{definition}
Thống kê Hotelling một mẫu là
\begin{equation}
 T^2=n(\Xbar-\vmu_0)'S^{-1}(\Xbar-\vmu_0).
 \label{eq:t2}
\end{equation}
\end{definition}
\nguon{Chương 5, §5.2, công thức (5-4), tr.~211}

\hieu{$T^2$ là bình phương khoảng cách từ trung bình mẫu đến vectơ
giả thuyết, đo bằng hiệp phương sai của \emph{trung bình mẫu}, ước lượng
bởi $S/n$. Giá trị lớn nghĩa là chênh lệch quan sát lớn so với sai số lấy mẫu.}

\begin{result}
Nếu $\mathbf X_1,\ldots,\mathbf X_n$ iid $\N_p(\vmu,\Sigma)$,
$\Sigma$ xác định dương, $n>p$, thì dưới $H_0$,
\begin{equation}
 F_{\mathrm{obs}}=\frac{n-p}{p(n-1)}T^2\sim F_{p,n-p}.
 \label{eq:t2F}
\end{equation}
Bác bỏ $H_0$ mức $\alpha$ nếu
\begin{equation}
 T^2>c_T^2,\qquad
 c_T^2=\frac{p(n-1)}{n-p}F_{p,n-p;1-\alpha}.
 \label{eq:t2critical}
\end{equation}
\end{result}
\nguon{Chương 5, §5.2, công thức (5-5)--(5-7), tr.~212}

Một cách tương đương là so sánh $F_{\mathrm{obs}}$ với phân vị $F$.
Không so trực tiếp $T^2$ với phân vị $F$ mà bỏ hệ số chuyển đổi.
Giá trị p là
\[
 \text{p-value}=1-G_{p,n-p}\!\left(\frac{n-p}{p(n-1)}T^2_{\mathrm{obs}}\right),
\]
trong đó $G_{p,n-p}$ là hàm phân phối tích lũy $F$.
Phần mềm thường có hàm survival function để tính đuôi phải ổn định hơn
phép trừ $1-G$ khi xác suất rất nhỏ.

\subsection{Ví dụ mẫu nhỏ: thực hiện trọn quy trình}
\begin{example}[Bài 3.1 của BTVN 2, tương ứng Exercise 5.1 trong sách]
Với dữ liệu bốn quan sát ở Phần~\ref{sec:sampling}, kiểm định
$\vmu_0=(7,11)'$ tại $\alpha=0.05$, giả sử mẫu iid chuẩn hai chiều.
\buoc{1}{Xác định kích thước và thống kê mẫu}
$n=4$, $p=2$, $\xbar=(6,10)'$, $S$ và $S^{-1}$ như
\eqref{eq:sample4}. Vì $|S|=44/9>0$, nghịch đảo tồn tại.
\buoc{2}{Tính vectơ chênh lệch}
$\mathbf d=\xbar-\vmu_0=(-1,-1)'$.
\buoc{3}{Tính thống kê}
\[
 T^2=\frac4{22}(-1,-1)
 \begin{pmatrix}9&15\\15&36\end{pmatrix}
 \begin{pmatrix}-1\\-1\end{pmatrix}
 =\frac{150}{11}\approx13.6364.
\]
\buoc{4}{Chọn phân phối và ngưỡng đúng}
\[
 F_{\mathrm{obs}}=\frac{4-2}{2(4-1)}T^2=
 \frac{T^2}{3}=\frac{50}{11}\approx4.5455.
\]
Dưới $H_0$, $F_{\mathrm{obs}}\sim F_{2,2}$.
Phân vị $F_{2,2;0.95}=19$, nên $c_T^2=3\cdot19=57$.
\buoc{5}{Ra quyết định và diễn giải}
$13.6364<57$: chưa bác bỏ $H_0$ ở mức 5\%.
Giá trị p bằng $11/61\approx0.1803$.
Chưa đủ bằng chứng để nói vectơ trung bình tổng thể khác $(7,11)'$.
\end{example}

Mẫu rất nhỏ làm $S$ có độ bất định lớn, nên ngưỡng kiểm định cao.
Nếu sai lầm dùng $\chi^2_{2;0.95}=5.9915$, ta sẽ bác bỏ ở ví dụ này,
trái với kiểm định chính xác. Đây không phải chi tiết nhỏ về làm tròn:
lựa chọn phân phối thay đổi kết luận.
\nguon{Dữ kiện Exercise 5.1, Chương 5; lý thuyết §5.2. Phép tính được kiểm tra độc lập}

\subsection{Kiểm tra bằng trường hợp một chiều}
Đặt $p=1$ trong \eqref{eq:t2F}:
\[
 \frac{n-1}{n-1}T^2=T^2\sim F_{1,n-1}.
\]
Mà $t_{n-1}^2\sim F_{1,n-1}$, nên kiểm định Hotelling thu về kiểm định
$t$ hai phía. Ngưỡng $F_{1,n-1;1-\alpha}$ bằng
$t_{n-1;1-\alpha/2}^2$.

\kiemtra{Trong ví dụ mẫu nhỏ, thống kê nào có phân phối $F_{2,2}$:
$T^2$ hay $T^2/3$? Đáp án: $T^2/3$.}

\subsection{Tính bất biến: đổi đơn vị đo không đổi kết luận}
\label{sec:invariance_proof}
Cho $\mathbf Y_j=C\mathbf X_j+\mathbf d$ với $C$ vuông khả nghịch.
Biến đổi đồng thời
\[
 \bar{\mathbf y}=C\xbar+\mathbf d,\quad S_Y=CS_XC',\quad
 \vmu_{Y,0}=C\vmu_0+\mathbf d.
\]
Khi đó
\begin{equation}
 T_Y^2=T_X^2.
 \label{eq:invariance}
\end{equation}
Đổi mét sang centimét, đổi thang nhiệt độ hoặc chọn một hệ trục khả nghịch
chỉ biểu diễn cùng thông tin theo cách khác. Chuẩn hóa chênh lệch bằng
hiệp phương sai làm các hệ số đổi đơn vị triệt tiêu.
\nguon{Chương 5, §5.2, phép biến đổi (5-9) và lập luận bất biến, tr.~215--216}

\buoc{1}{Chứng minh các thống kê mẫu biến đổi đồng bộ}
Tính trung bình hai vế của $\mathbf y_j=C\mathbf x_j+\mathbf d$ cho
$\bar{\mathbf y}=C\xbar+\mathbf d$. Do đó
\[
 \mathbf y_j-\bar{\mathbf y}=C(\mathbf x_j-\xbar).
\]
Thế vào định nghĩa hiệp phương sai mẫu:
\begin{align*}
 S_Y&=\frac1{n-1}\sum_j
 C(\mathbf x_j-\xbar)(\mathbf x_j-\xbar)'C'\\
 &=C\left[\frac1{n-1}\sum_j
 (\mathbf x_j-\xbar)(\mathbf x_j-\xbar)'\right]C'=CS_XC'.
\end{align*}
Giả thuyết cũng đổi hệ tọa độ:
$\vmu_{Y,0}=C\vmu_0+\mathbf d$.
Đặt $\mathbf r=\xbar-\vmu_0$, ta nhận được độ lệch mới
\[
 \bar{\mathbf y}-\vmu_{Y,0}=C\mathbf r.
\]
Phần tịnh tiến triệt tiêu vì xuất hiện ở cả hai trung bình.

\buoc{2}{Thiết lập nghịch đảo của hiệp phương sai mới}
Với tích các ma trận khả nghịch, nghịch đảo lấy theo thứ tự ngược lại:
\[
 (CS_XC')^{-1}=(C')^{-1}S_X^{-1}C^{-1}.
\]
Có thể kiểm tra bằng cách nhân với $CS_XC'$: các cặp
$C'(C')^{-1}$, $S_XS_X^{-1}$, $CC^{-1}$ lần lượt triệt tiêu và cho $I_p$.

\buoc{3}{Thế vào thống kê Hotelling}
\begin{align*}
 T_Y^2
 &=n(C\mathbf r)'(CS_XC')^{-1}(C\mathbf r)\\
 &=n\mathbf r'C'(C')^{-1}S_X^{-1}C^{-1}C\mathbf r\\
 &=n\mathbf r'S_X^{-1}\mathbf r=T_X^2.
\end{align*}
Các hệ số đổi tọa độ xuất hiện trong cả độ lệch lẫn hiệp phương sai,
nên triệt tiêu khi chuẩn hóa. Đây là chứng minh bất biến của \emph{giá
trị thống kê}, không chỉ của dạng công thức.

\buoc{4}{Suy ra bất biến của giá trị p và quyết định kiểm định}
$C$ là ma trận $p\times p$ khả nghịch nên không thay đổi $n$ hoặc $p$.
Hệ số chuyển sang $F$, số bậc tự do và ngưỡng $c_T^2$ đều giữ nguyên.
Vì vậy cùng một $T^2$ cho cùng giá trị p và cùng quyết định bác bỏ hay
chưa bác bỏ. Dưới giả thiết mẫu chuẩn, mẫu biến đổi vẫn chuẩn và có
hiệp phương sai $C\Sigma C'\succ0$, nên phân phối kiểm định chính xác
vẫn áp dụng được.

\hieu{Nếu đổi đơn vị đo hoặc trộn các biến theo một cách có thể hoàn
tác, ta vẫn giữ toàn bộ thông tin về giả thuyết. Kiểm định vectơ trung
bình cần cho cùng kết luận trong hai cách biểu diễn đó.}

Đẳng thức đại số $T_Y^2=T_X^2$ không cần giả thiết chuẩn, chỉ cần các
nghịch đảo tồn tại và phép biến đổi phù hợp. Giả thiết chuẩn cần để dùng
phân phối $F$ \emph{chính xác}; không nên nhập hai loại điều kiện làm một.

\begin{example}[Bất biến từ số liệu của Example 5.1]
Dữ liệu $(6,9),(10,6),(8,3)$ cho
$\xbar=(8,6)'$, $S_X=\begin{pmatrix}4&-3\\-3&9\end{pmatrix}$.
Với $\vmu_0=(9,5)'$, ta có $T_X^2=7/9$.
Chọn $C=\begin{pmatrix}1&-1\\1&1\end{pmatrix}$, $\mathbf d=0$:
\[
 \bar{\mathbf y}=(2,14)',\quad\vmu_{Y,0}=(4,14)',\quad
 S_Y=\begin{pmatrix}19&-5\\-5&7\end{pmatrix}.
\]
$S_Y^{-1}=\frac1{108}\begin{pmatrix}7&5\\5&19\end{pmatrix}$, nên
\[
 T_Y^2=\frac3{108}(-2,0)
 \begin{pmatrix}7&5\\5&19\end{pmatrix}(-2,0)'=\frac79.
\]
\end{example}
\nguon{Chương 5, Example 5.1, tr.~213; biến đổi dùng trong Bài 3.2 của BTVN 2}

\luuy{Phải biến đổi cả giả thuyết. Nếu đổi dữ liệu mà giữ nguyên
$\vmu_0$, ta kiểm định một giả thuyết khác. Nếu $C$ không khả nghịch,
phép biến đổi có thể làm mất thông tin và tính bất biến này không còn áp dụng.}

Trong hình học, miền tin cậy của dữ liệu biến đổi chính là ảnh
$\{C\vmu+\mathbf d:\vmu\in\mathcal C_X\}$ của miền ban đầu.
Elipsoid có thể xoay, co giãn hoặc bị kéo xiên, nhưng một vectơ giả
thuyết và ảnh của nó luôn có cùng trạng thái nằm trong/ngoài miền.
Đây là cách nhìn hình học của cùng chứng minh.
Chọn một tập con biến, tương ứng với ma trận chữ nhật làm giảm số
chiều, không phải phép đổi tọa độ khả nghịch và không thuộc kết quả này.

\subsection{Tỷ số hợp lý: so hai cách giải thích cùng dữ liệu}
Hàm hợp lý $L(\vmu,\Sigma)$ cho biết mức độ phù hợp của tham số với
dữ liệu đã quan sát. Ta so sánh:
\begin{itemize}
 \item mô hình không ràng buộc: được chọn cả $\vmu$ và $\Sigma$;
 \item mô hình dưới $H_0$: giữ $\vmu=\vmu_0$, chỉ được chọn $\Sigma$.
\end{itemize}
Mô hình có ràng buộc không thể đạt cực đại lớn hơn mô hình tự do. Tỷ số
giữa hai cực đại nằm trong $[0,1]$; tỷ số nhỏ nghĩa là ràng buộc làm giảm
độ phù hợp nhiều.

Đặt $\mathbf d=\xbar-\vmu_0$. Các MLE là
\begin{align}
 \widehat\Sigma&=\frac{n-1}{n}S,\\
 \widehat\Sigma_0&=\frac1n\sum_j(\mathbf x_j-\vmu_0)(\mathbf x_j-\vmu_0)'
 =\widehat\Sigma+\mathbf d\mathbf d'.
\end{align}
\nguon{Chương 5, §5.3, công thức (5-10)--(5-12), tr.~216--217}

Để tránh nhầm tên, bài giảng dùng hai ký hiệu riêng:
\begin{align}
 \Lambda_{\mathrm{LR}}
 &=\frac{\max_{\Sigma}L(\vmu_0,\Sigma)}{\max_{\vmu,\Sigma}L(\vmu,\Sigma)}
 =\left(\frac{|\widehat\Sigma|}{|\widehat\Sigma_0|}\right)^{n/2},\\
 \lambda_W&=\Lambda_{\mathrm{LR}}^{2/n}
 =\frac{|\widehat\Sigma|}{|\widehat\Sigma_0|}
 =\left(1+\frac{T^2}{n-1}\right)^{-1}.
 \label{eq:wilks}
\end{align}
$\lambda_W$ là Wilks' lambda trong bài toán này.
Trong sách, $\Lambda$ chỉ tỷ số hợp lý và $\Lambda^{2/n}$ là Wilks' lambda.
Từ \eqref{eq:wilks},
\begin{equation}
 T^2=(n-1)\left(\frac{|\widehat\Sigma_0|}{|\widehat\Sigma|}-1\right).
 \label{eq:t2det}
\end{equation}
\nguon{Chương 5, §5.3, Result 5.1, công thức (5-14), tr.~218; (5-15), tr.~219}

\hieu{$T^2$ lớn tương đương $\lambda_W$ nhỏ. Hai cách xây dựng dẫn
đến cùng kiểm định trong bài toán chuẩn một mẫu. Điều này nối ý tưởng
khoảng cách với ý tưởng so độ phù hợp mô hình.}

\begin{example}[Kiểm tra lại bằng định thức]
Ở mẫu $n=4$ trên, đặt $A=3S$, $A_0=A+4\mathbf d\mathbf d'$.
Ta có
\[
 A=\begin{pmatrix}24&-10\\-10&6\end{pmatrix},\qquad
 A_0=\begin{pmatrix}28&-6\\-6&10\end{pmatrix},
\]
nên $|A|=44$, $|A_0|=244$.
Hệ số chia $n$ trong hai MLE triệt tiêu khi lấy tỷ số định thức:
\[
 T^2=3\left(\frac{244}{44}-1\right)=\frac{150}{11}.
\]
$\lambda_W=11/61$ và $\Lambda_{\mathrm{LR}}=(11/61)^2$.
Việc $\lambda_W$ bằng giá trị p trong ví dụ này là trùng hợp riêng của
$n=4,p=2$, không phải công thức tổng quát cho giá trị p.
\end{example}

Trong tính toán lớn, log-định thức thường ổn định hơn định thức trực tiếp.
Tuy vậy cách tính ưu tiên cho $T^2$ vẫn là giải hệ với $S$; công thức định
thức ở đây giúp hiểu lý thuyết và kiểm tra chéo.
